%
% Elementare Funktionen
%
\subsection{Elementare Funktionen und das Integrationsproblem}
Der Prozess der Erweiterung des Differentialkörpers der 
rationalen Funktionen $\mathbb{Q}(x)$ durch algebraische Elemente,
Logarithmen und Exponentialfunktionen reproduziert also die Vorgehensweise
eines Mathematikers, der ein Integralproblem lösen will.
Er beginnt mit der Konstruktion eines Funktionenkörpers, der den
Integranden enthält.
Dazu wird es notwendig sein, zusätzliche Konstanten wie Wurzeln,
$\pi$, $i$ und andere oder Parameter hinzuzufügen.
Ausserdem müssen möglicherweise wiederholt Exponentialfunktionen,
Logarithmen und algebraische Funktionen hinzugefügt werden.
Alle Funktionen, mit denen ein menschlicher Rechner arbeitet, sind
auf diese Art konstruierbar.
Dies rechtfertigt, diesen Funktionen einen eigenen Namen zu geben.

\begin{definition}[Elementare Funktionen]
Sei $D$ eine Derivation in $\mathbb{Q}(x)$ mit $Dx=1$.
Eine Funktion $f(x)$ heisst \emph{elementar}, wenn $f$ in einer
Erweiterung von $\mathbb{Q}(x)$ enthalten ist, die durch hinzufügen
algebraischer Elemente, Logarithmen oder Exponentialfunktionen
entsteht.
\index{elementare Funktion}%
\end{definition}

Mit dem Begriff der elementaren Funktion lässt sich jetzt das
Problem der Integration in geschlossener Form präzis formulieren.
\index{geschlossene Form}%
Ein menschlicher Rechner erwartet, dass er den Differentialkörper,
in dem der Integrand gefunden werden kann, durch weitere Erweiterungen
so vergrössern kann, dass sich auch die Stammfunktion darin finden
lässt.
Die Stammfunktion ist daher auch eine elementare Funktion.
Das Problem der Integration in geschlossener Form ist daher das
Folgende.

\begin{problem}[Integration in geschlossener Form]
\index{Integration in geschlossener Form}%
Gegeben eine elementare Funktion $f$ in einem Differentialkörper
$F$, finde eine elementare Funktion $g$ mit $Dg = f$.
\end{problem}
